\section{Experimental details and Additional Results}\label{appendix experimental details}

\subsection{Datasets and featurization}\label{apdx datasets and fingerprints}

The 1000 molecules from GuacaMol are included in our code.
The molecules from the \textsc{dockstring} dataset are available online (\url{https://github.com/dockstring/dataset}).

Note that when using count fingerprints for $T_{DP}$ we used the \emph{square root} of the counts as the fingerprint.
This was done for two reasons:
\begin{enumerate}
    \item To reduce the norms of the vectors (and thereby increase $\zeta$).
    \item To roughly make their interpretation consistent with $T_{MM}$: i.e.\@ the ``weight'' of a fragment in which occurs $n$ times in the numerator/denominator of $T_{DP}$ is $n$ times the weight of a fragment which occurs once.
\end{enumerate}

\subsection{Details of $T_{MM}$ random features}\label{apdx:explain tmm features and hash}

We use the random hash from \citet{ioffe2010improved} in our implementation.
To hash vectors in $\mathbb{R}^D$
random variables
$r_i, c_i\sim \Gamma(2, 1)$, $\beta_i\sim U(0, 1)$
are drawn i.i.d.\@ for $i=1,\ldots,D$,
and then the following\footnote{
Note that the presentation of equations~\ref{eqn:ioffe-hash-first}--\ref{eqn:ioffe-hash-last}
differs slightly from \citet{ioffe2010improved} who defines $y_i$ and $a_i$ to be the exponential
of equations~\ref{eqn:ioffe-hash-second}/\ref{eqn:ioffe-hash-third}: we wrote it this way because in practice
working in log space avoids numerical stability issues.
This is explicitly suggested in their paper.
}
are computed for all $i$:
\begin{align}
    t_i(x) &= \lfloor (\ln{x_i})/r_i + \beta_i \rfloor \label{eqn:ioffe-hash-first}\,, \\
    y_i(x) &= r_i(t_i(x)-\beta_i) \label{eqn:ioffe-hash-second}\,, \\
    a_i(x) &= \ln{c_i} - y_i(x) - \ln{r_i} \label{eqn:ioffe-hash-third}\,, \\
    i^*(x) &= \arg\min_{i=1,\ldots,D} a_i(x)\,,  \\
    h_{r,c,\beta}(x) &= \left(i^*(x), t_{i^*(x)}(x)\right) \,.\label{eqn:ioffe-hash-last}
\end{align}
Note that equation~\ref{eqn:ioffe-hash-first} uses the convention that $\ln{0}=-\infty$.
These variables do not have a clear interpretation in isolation so we refer the reader
to \citet{ioffe2010improved} for an explanation of why this hashing procedure
produces a random hash for $T_{MM}$.
The hash itself is formed of 2 integers: $i^*\in\{1,\ldots,D\}$
and $t_{i^*}\in\mathbb{Z}$.
This unfortunately means that the number of possible outputs is potentially infinite,
which would require us to potentially sample an arbitrarily large vector $\Xi$.
To avoid this, we first use python's built-in \texttt{hash} function to map this pair of integers to a single integer,
then take the result module $2^{12}=4096$.
This allows us to sample a small (finite) vector $\Xi$,
and although it introduces a small amount of bias this does not appear to be an issue in practice.

Elements of $\Xi$ are always sampled i.i.d.\@ from either a Rademacher or Gaussian distribution
(the distribution should always be clear from context).


\subsection{Additional results for $T_{DP}$ random features}

\begin{figure}
    \centering
    \includegraphics{new-figures/poly_and_term_term_mse.pdf}
    \caption{
        MSEs for approximating $(x\cdot x')^r$ (\textbf{left})
        and $((x\cdot x')/(\|x\|^2+\|x'\|^2))^r$ (\textbf{right})
        using \textsc{TensorSketch} for various $r$ on count and binary fingerprints
        as a function of the random feature dimension $M$.
    }
    \label{fig:tdp-polynomial-and-single-term-error}
\end{figure}

\begin{figure}
    \centering
    \includegraphics{new-figures/tdp_feat_alloc_and_overall_error_count.pdf}
    \caption{
    Same as Figure~\ref{fig:tdp-total-rf} but with count fingerprints.
    }
    \label{fig:tdp-total-rf-count-apdx}
\end{figure}

Figure~\ref{fig:tdp-polynomial-and-single-term-error} shows that the error for both polynomial random features
from \textsc{TensorSketch} and random features for individual terms in $T_{DP}'s$ power series
decreases with $O(1/M)$.
Figure~\ref{fig:tdp-total-rf-count-apdx} shows the overall error for count fingerprints.



\subsection{Gaussian process training details from section~\ref{sec:regression-expt}}\label{apdx:regression-expt-details}

Our GP models use a constant mean and Gaussian noise.
Specifically, the model of the observed labels $y$ is:
\begin{align*}
    f(X) &\sim \mathcal{GP}\left(\mu, ak(X,X)]\right) \\
    y(X) &\sim \mathcal{N}\left(f(X), \sigma^2I\right)
\end{align*}
Therefore, the GP hyperparameters are three scalars:
\begin{itemize}
    \item The constant mean, $\mu$
    \item The kernel amplitude/outputscale $a$
    \item The observation noise $\sigma^2$
\end{itemize}

GP performance will be greatly affected by the choice of kernel hyperparameters.
To ensure that the difference in performance is not due to differences in kernel hyperparameter settings
we fit them in a consistent way for all methods.
Specifically,
for all methods,
we start by fitting an exact GP to a random subset of $M=5000$ data points by maximizing the marginal likelihood with L-BFGS.
The different approximations are as follows:
\begin{itemize}
    \item \textbf{Random subset:}
        \texttt{gpytorch}'s exact GP inference is applied on a random subset of $M$ data points
        (a different subset then was used to fit the hyperparameters).
    \item \textbf{SVGP:}
        \texttt{scikit-learn}'s K-means clustering is run with $M$ clusters to produce an initialization of the inducing points.
        These inducing points are used to initialize a sparse variational GP \citep{hensman2013gaussian},
        implemented in \texttt{gpytorch}.
        The variational parameters are optimized via natural gradient descent with a learning rate of $10^{-1}$
        and a batch size of $2M=10\,000$ for one pass through the dataset.
        Although the inducing point locations themselves could be further optimized with gradient descent,
        we chose not to do so for this experiment.
    \item \textbf{RFGP:}
        First, the training and test sets are converted into $M$ dimensional random features.
        Then, the posterior equations for inference in Bayesian linear models
        are used to make predictions on the test set given the training set \citep[equations 3.49--3.51]{bishop2006pattern}.
        The computation is done in a specific order to avoid forming any $n\times n$ matrices
        (the Woodbury matrix identity is used extensively for this).
        This is fairly clearly documented in the code.
\end{itemize}

Note that for $T_{MM}$ it was vital to implement the kernel as
\begin{equation}\label{eqn:minmax-using-l1-distance}
    T_{MM}(x, x') = \frac{\|x\|_1+\|x'\|_1 - \|x-x'\|_1}{\|x\|_1+\|x'\|_1 + \|x-x'\|_1}
\end{equation}
instead of a naive implementation which follows equation~\ref{eqn:tanimoto-for-sets}.
This is because such an implementation requires forming a tensor of shape $N\times M \times d$
when calculating the kernel between $N$ and $M$ points in $\mathbb{R}^d$
(for example $T_{ijk}=\min(x_{ik}, y_{jk})$)
which can exceed memory limits for modest $N,M,d$.
This identity allows us to use the relatively efficient \texttt{torch.cdist} function.
Note that we did \emph{not} invent this identity ourselves;
we discovered it in \citet{ioffe2010improved}.

Even with this implementation, $M=5000$ inducing points did not fit into GPU memory for $T_{MM}$,
so all experiments were run on CPU.

\subsection{Metrics and additional results from section~\ref{sec:regression-expt}}\label{apdx:regression-expt-extra-results}

\paragraph{Metrics}
Table~\ref{tab:gp-regression-logp} reports \emph{average log probability},
which calculated by first calculating the log probability of each test point individually
(i.e.\@ marginally, not jointly), then averaging these values.
We also report the coefficient of determination ($R^2$ score),
calculated using the function \texttt{sklearn.metrics.r2\_score}.
This measures only the error of the GP mean.
A value of 1 indicates perfect prediction,
while a value of 0 can be achieved by predicting the sample mean for every data point.

\paragraph{Additional results}
Table~\ref{tab:gp-regression-r2} reports the average $R^2$.
Trends are similar to Table~\ref{tab:gp-regression-logp}.
We also include baselines for two types of graph neural network:
Attentive FP \citep{xiong2019pushing}
and MPNN \citep{gilmer2017neural}.
Although the performance of GP methods does not match that of Attentive FP,
it is often close,
suggesting there is potential for approximate GPs to be competitive with graph neural networks
for molecular property prediction.

\begin{table*}[tb]
\caption{
Average $R^2$ score for approximate GPs on \textsc{dockstring} dataset.
Attentive FP and MPNN results
are taken from \citet{garcia2022dockstring}.
Other details are the same as Table~\ref{tab:gp-regression-logp}.
}
\label{tab:gp-regression-r2}
\begin{center}
\resizebox{0.95\textwidth}{!}{
\begin{sc}
\begin{tabular}
{llr@{\hspace{0.02cm}$\pm$\hspace{0.02cm}}l@{\hspace{0.30cm}}r@{\hspace{0.02cm}$\pm$\hspace{0.02cm}}l@{\hspace{0.30cm}}r@{\hspace{0.02cm}$\pm$\hspace{0.02cm}}l@{\hspace{0.30cm}}r@{\hspace{0.02cm}$\pm$\hspace{0.02cm}}l@{\hspace{0.30cm}}r@{\hspace{0.02cm}$\pm$\hspace{0.02cm}}l@{\hspace{0.30cm}}}
\toprule
Kernel & Method & \multicolumn{2}{c}{ESR2} & \multicolumn{2}{c}{F2} & \multicolumn{2}{c}{KIT} & \multicolumn{2}{c}{PARP1} & \multicolumn{2}{c}{PGR}\\
\midrule
$T_{MM}$ & Rand subset GP & 0.514 & 0.002 & 0.810 & 0.002 & 0.695 & 0.002 & 0.849 & 0.001 & 0.426 & 0.007\\
 & SVGP & 0.578 & 0.001 & 0.861 & 0.000 & 0.749 & 0.000 & 0.889 & 0.000 & 0.542 & 0.002\\
 & RFGP ($\Xi$ Rad.) & 0.518 & 0.002 & 0.838 & 0.001 & 0.703 & 0.002 & 0.864 & 0.001 & 0.465 & 0.003\\
 & RFGP ($\Xi$ Gauss.) & 0.517 & 0.002 & 0.837 & 0.000 & 0.702 & 0.001 & 0.864 & 0.001 & 0.467 & 0.004\\
\midrule
$T_{DP}$ & Rand subset GP & 0.513 & 0.003 & 0.817 & 0.001 & 0.696 & 0.002 & 0.851 & 0.001 & 0.384 & 0.011\\
 & SVGP & 0.581 & 0.001 & 0.865 & 0.000 & 0.753 & 0.001 & 0.889 & 0.000 & 0.543 & 0.002\\
 & RFGP (plain) & 0.546 & 0.001 & 0.852 & 0.001 & 0.716 & 0.002 & 0.876 & 0.000 & 0.512 & 0.002\\
 & RFGP (norm) & 0.546 & 0.001 & 0.852 & 0.001 & 0.715 & 0.002 & 0.876 & 0.000 & 0.513 & 0.002\\
 & RFGP (sketch) & 0.545 & 0.001 & 0.852 & 0.001 & 0.716 & 0.002 & 0.876 & 0.000 & 0.510 & 0.002\\
\midrule
N/A & MPNN & 0.506 & 0.001  & 0.798 & 0.005  & 0.755 & 0.005  & 0.815 &  0.010 & 0.324 &  0.096 \\
N/A & Attentive FP & 0.627 &  0.010 & 0.880 &  0.001 & 0.806 &  0.008 & 0.910 &  0.002 & 0.678 & 0.008 \\
\bottomrule
\end{tabular}
\end{sc}
}
\end{center}
\end{table*}
